\section{Materials and methods}
\label{sec:methods}

Fig.~\ref{fig:pipeline} summarises the framework. This study is reported following the TRIPOD+AI recommendations where applicable~\cite{collins2024tripodai}.

\begin{figure*}[!htbp]
  \centering
  \includegraphics[width=\textwidth]{fig1_pipeline.pdf}
  \caption{Proposed pipeline. Imputation, feature engineering, standardisation and RFECV are fitted on the training part of each split only. The features retained by RFECV are passed to eight heterogeneous base learners, whose out-of-fold (OOF) probabilities train a logistic-regression meta-learner.}
  \label{fig:pipeline}
\end{figure*}

\subsection{Data}
We used the PIMA Indians Diabetes dataset~\cite{smith1988adap}: 768 women of Pima heritage aged at least 21 years, of whom 268 (34.9\%) had diabetes. Each record contains eight predictors: number of pregnancies, 2-h plasma glucose concentration in an oral glucose tolerance test (OGTT; mg/dL), diastolic blood pressure (mm\,Hg), triceps skin-fold thickness (mm), 2-h serum insulin (\textmu U/mL), BMI (kg/m$^2$), diabetes pedigree function and age (years). The data are publicly available and de-identified~\cite{pimadata}.

\paragraph{Outcome definition and intended setting} According to the dataset documentation, the binary outcome indicates whether the woman showed signs of diabetes according to World Health Organization criteria, i.e.\ a 2-h post-load plasma glucose of at least 200~mg/dL at a survey examination or diagnosis during routine medical care~\cite{pimadata,smith1988adap}. This outcome is not a contemporary clinical diagnosis of T2DM based on current HbA1c or fasting-glucose criteria~\cite{ada2024}. Moreover, the strongest predictor, 2-h OGTT glucose, is itself a diagnostic measurement closely related to the outcome definition. The models studied here therefore address risk stratification among individuals for whom OGTT measurements are already available, not low-cost population screening. We treat the analysis primarily as a methodological benchmark rather than as the development of a clinically deployable diagnostic model.

\subsection{Missing values}
Zero is physiologically impossible for glucose, blood pressure, skin-fold thickness, insulin and BMI, so these zeros were recoded as missing (Table~\ref{tab:missing-counts}). In the primary analysis, missing values were replaced by the median of the training data in each split. Because insulin is missing for almost half of the patients, we also compared three alternatives in a sensitivity analysis (Section~\ref{sec:missing-methods}).

\begin{table}[!htbp]
\centering
\caption{Physiologically impossible zero values recoded as missing ($n = 768$).}
\label{tab:missing-counts}
\begin{tabular}{lrr}
\toprule
Predictor & Missing ($n$) & Missing (\%) \\
\midrule
Insulin & 374 & 48.7 \\
Skin-fold thickness & 227 & 29.6 \\
Diastolic blood pressure & 35 & 4.6 \\
BMI & 11 & 1.4 \\
Glucose & 5 & 0.7 \\
\bottomrule
\end{tabular}
\end{table}

\subsection{Feature engineering}
Eight derived features encoding known T2DM risk interactions were added after imputation, giving 16 candidate predictors for RFECV: the 8 original variables and 8 deterministic derived features (Table~\ref{tab:features}). Because the derived features are functions of the original variables, the 16 candidates are not informationally independent. The insulin-resistance term follows the HOMA-IR formula~\cite{matthews1985homa}, but PIMA records 2-h post-load rather than fasting glucose and insulin. We therefore name it \texttt{HOMA\_IR\_surrogate} and do not interpret it as clinical HOMA-IR. The binary thresholds follow the WHO obesity cut-off~\cite{who2000obesity} and the 2-h OGTT threshold for impaired glucose tolerance~\cite{ada2024}.

\begin{table}[!htbp]
\centering
\caption{Derived features. All continuous inputs are measured after imputation.}
\label{tab:features}
\small
\begin{tabularx}{\linewidth}{llL}
\toprule
Feature & Definition & Rationale \\
\midrule
\texttt{Glucose\_BMI} & Glucose $\times$ BMI / 100 & Joint glycaemia--adiposity effect \\
\texttt{Glucose\_Age} & Glucose $\times$ Age / 100 & Glycaemic risk rising with age \\
\texttt{HOMA\_IR\_surrogate} & Glucose $\times$ Insulin / 405 & Insulin-resistance surrogate (post-load) \\
\texttt{Insulin\_Glucose\_Ratio} & Insulin / Glucose & Insulin response relative to glucose \\
\texttt{BMI\_Age} & BMI $\times$ Age / 100 & Cumulative adiposity exposure \\
\texttt{Obese} & $1$ if BMI $\geq 30$ & WHO obesity threshold \\
\texttt{IGT\_2h} & $1$ if Glucose $\geq 140$ & Impaired glucose tolerance threshold \\
\texttt{Pedigree\_Age} & Pedigree $\times$ Age & Familial risk expressed over time \\
\bottomrule
\end{tabularx}
\end{table}

\subsection{Recursive feature elimination with cross-validation}
All features were standardised with training statistics. RFECV started from the 16 candidate predictors. At each step it fitted a random forest (200 trees, balanced class weights), ranked features by impurity-based importance and removed the least important one. Each subset size was scored by stratified 5-fold cross-validated ROC-AUC within the training data, and the size with the highest mean score was retained (minimum three features).

\paragraph{Common versus learner-specific selection}
In the model comparison, RFECV was fitted once within each outer training fold using the random-forest ranking estimator, and the resulting subset was supplied to every candidate learner. This gives all learners an identical, leakage-free feature-selection protocol at the cost of one RFECV fit per fold. The subset is optimised with the random-forest ranking criterion, not separately for each learner. It is therefore \emph{not} equivalent to cross-validating a learner-specific RFECV pipeline for every classifier. To gauge the effect of this choice, we also evaluated logistic regression with its own RFECV ranked by absolute standardised coefficients (``LR-ranked RFECV''). Throughout the article, ``+ RFECV'' denotes the common random-forest-ranked RFECV subset.

\subsection{Stacking ensemble}
The level-0 learners were chosen for diversity of inductive bias (Table~\ref{tab:learners}). The level-1 meta-learner was an L2-regularised logistic regression with balanced class weights. It was trained on the positive-class probabilities produced by each base learner for held-out samples in an internal stratified 5-fold split~\cite{wolpert1992stacked,ting1999issues}. The base learners were then refitted on all training data for prediction. Gaussian naive Bayes and a depth-5 decision tree were evaluated as additional reference learners but were not part of the stack. Hyperparameters were fixed a priori at conventional settings for small tabular data (Table~\ref{tab:learners}): library defaults where these are standard (e.g.\ $C=1$ for LR and SVM), moderate ensemble sizes (200--300 trees), shallow boosting (depth 3) with a small learning rate (0.05), and a neighbourhood of $k=15$. They were not tuned, so that no tuning loop could leak into the evaluation. All comparisons between learners are therefore conditional on these prespecified configurations and do not cover every possible tuned version of each algorithm.

\begin{table}[!htbp]
\centering
\caption{Level-0 learners and their fixed settings.}
\label{tab:learners}
\small
\begin{tabularx}{\linewidth}{llL}
\toprule
Learner & Family & Settings \\
\midrule
Logistic regression (LR) & Linear & L2, $C=1$, balanced class weights \\
SVM & Kernel & RBF, $C=1$, balanced weights, Platt-calibrated~\cite{platt1999} \\
$k$-nearest neighbours & Instance-based & $k=15$, distance-weighted \\
Random forest~\cite{breiman2001rf} & Bagged trees & 300 trees, min.\ leaf 2, balanced weights \\
Extra trees & Randomised trees & 300 trees, min.\ leaf 2, balanced weights \\
Gradient boosting & Boosted trees & 200 trees, depth 3, learning rate 0.05 \\
XGBoost~\cite{chen2016xgboost} & Boosted trees & 300 trees, depth 3, rate 0.05, subsample 0.8 \\
LightGBM~\cite{ke2017lightgbm} & Boosted trees & 300 trees, 15 leaves, rate 0.05, balanced weights \\
\bottomrule
\end{tabularx}
\end{table}

\subsection{Internal evaluation design}
Supplementary Fig.~S1 shows where each loop runs. The data were split once into a stratified 80\% development set (614 patients, 214 with diabetes) and a 20\% hold-out set (154 patients, 54 with diabetes), with random seed 42. The split was fixed in code before any model was fitted. Hold-out results of the same prespecified pipeline were inspected in an earlier draft of this work; no learner, hyperparameter, feature or preprocessing choice was changed afterwards, and the hold-out set was not used for any modelling decision.

\paragraph{Repeated cross-validation}
On the development set, 11 learners (the eight base learners, naive Bayes, the decision tree and the stack) were compared with and without the common RFECV subset (22 configurations), together with logistic regression with LR-ranked RFECV (23 configurations in total; Supplementary Table~S5), using stratified 10-fold cross-validation repeated five times (50 folds). Within each fold, imputation, feature engineering, standardisation and RFECV were fitted on the nine training folds only. Because fold estimates from overlapping training sets are correlated, we report the mean of each metric with a 95\% confidence interval based on the Nadeau--Bengio corrected variance~\cite{nadeau2003inference},
\begin{equation}
  \widehat{\operatorname{Var}}_{\mathrm{corr}} = \left(\frac{1}{J} + \frac{n_{\mathrm{test}}}{n_{\mathrm{train}}}\right) s^2,
  \label{eq:corrected}
\end{equation}
where $J=50$ is the number of fold estimates, $s^2$ their sample variance, and $n_{\mathrm{test}}/n_{\mathrm{train}} = 1/9$. Each configuration was compared with logistic regression on all features using the corrected resampled $t$-test on paired fold differences ($J-1$ degrees of freedom)~\cite{nadeau2003inference,bouckaert2004replicability}. $p$-values were adjusted for the 22 comparisons with the Holm procedure~\cite{holm1979}. The comparison set was fixed in the analysis code before the repeated cross-validation was run. We chose this test because it is the standard, simulation-validated correction for the dependence between repeated cross-validation folds and has acceptable type~I error and good replicability~\cite{nadeau2003inference,bouckaert2004replicability}. A nested or paired bootstrap of the whole pipeline would have required refitting RFECV and the stack thousands of times (about 45~s per fit). The full testing procedure is given in Supplementary Section~S2.

\paragraph{Hold-out evaluation}
The final pipelines were fitted on the full development set and scored once on the hold-out set. ROC-AUC confidence intervals were computed with DeLong's method, and ROC-AUCs of correlated models were compared with DeLong's test~\cite{delong1988,sun2014delong}, with Holm adjustment over the six comparisons. All other metrics were given 95\% percentile bootstrap confidence intervals (2000 resamples of the hold-out patients).

\subsection{Performance measures}
We report accuracy, balanced accuracy, precision (positive predictive value, PPV), recall (sensitivity), specificity, F1, the Matthews correlation coefficient (MCC)~\cite{chicco2020mcc}, ROC-AUC and average precision. Calibration was assessed with the Brier score, calibration-in-the-large (the intercept $a$ of a logistic model with $\operatorname{logit}\hat p$ as offset), the calibration slope $b$ from $\operatorname{logit} P(y=1) = a' + b\,\operatorname{logit}\hat p$, and calibration curves~\cite{vancalster2019calibration}. A well-calibrated model has $a = 0$ and $b = 1$. Calibration was evaluated both on the pooled out-of-fold predictions of the development set and on the hold-out set. The pooled set contains 3070 predictions: each of the 614 unique development patients is predicted once in each of the five repeats. No recalibration was applied to the reported models. For reference, the development-set prevalence is $\pi = 214/614 = 0.349$, so an intercept correction for balanced class weighting would be $\operatorname{logit}(\pi) = -0.626$.

\subsection{Threshold and decision-curve analysis}
Instead of using the default 0.5 threshold, we pre-specified operating thresholds targeting 80\%, 85\%, 90\% and 95\% sensitivity. Each threshold was chosen on the development-set out-of-fold predictions, using one probability per patient (the mean over the five repeats). The threshold for a target sensitivity $s$ is the $\lceil s\,n_+\rceil$-th largest predicted probability among the $n_+ = 214$ development positives. Patients are flagged when their probability is greater than or equal to the threshold, so ties at the threshold are flagged and the development sensitivity is at least the target. The thresholds were then applied unchanged to the hold-out set, where sensitivity, specificity, PPV, NPV and the proportion of patients flagged were estimated with bootstrap confidence intervals. Potential clinical utility was assessed by internal decision-curve analysis~\cite{vickers2006dca} on the hold-out set, using the net benefit
\begin{equation}
  \mathrm{NB}(p_t) = \frac{\mathrm{TP}(p_t)}{n} - \frac{\mathrm{FP}(p_t)}{n}\cdot\frac{p_t}{1-p_t},
\end{equation}
over threshold probabilities $p_t$ from 0.05 to 0.60, and compared with treating all or no patients.

\subsection{Missing-data sensitivity analysis}
\label{sec:missing-methods}
We repeated the development-set comparison of logistic regression and stacking + RFECV (10-fold cross-validation repeated twice) with four imputation strategies, each fitted within the training folds: (i) median imputation; (ii) median imputation plus a binary missingness indicator for every predictor with missing values, which lets the model use informative missingness~\cite{sperrin2020missing}; (iii) iterative chained-equation imputation (scikit-learn \texttt{IterativeImputer}, Bayesian ridge, 20 iterations)~\cite{vanbuuren2011mice}; and (iv) $k$-nearest-neighbour imputation ($k=10$)~\cite{troyanskaya2001knn}. Differences from median imputation were tested with the corrected resampled $t$-test.

\subsection{Explainability}
We computed Shapley values~\cite{lundberg2017shap} for the hold-out predictions of the final stacking + RFECV pipeline. The pipeline was treated as one function of the eight raw clinical inputs. Because there are only eight inputs, the values were obtained by exact enumeration of all $2^8 = 256$ feature coalitions (SHAP \texttt{ExactExplainer}), with features outside a coalition replaced by values from 40 development-set patients drawn by simple random sampling with seed 42 (an interventional background). They are therefore exact with respect to this background distribution. SHAP values describe how much each input contributed to the model's predictions; they are attributions, not estimates of causal effects. Explanations are therefore expressed in clinically measured variables, not in engineered or standardised features.

\subsection{Implementation and reproducibility}
The framework was implemented in Python 3.11 with scikit-learn 1.9~\cite{pedregosa2011sklearn}, XGBoost 3.2, LightGBM 4.7 and SHAP. Every stochastic step used seed 42: the hold-out split, the repeated outer folds, the inner RFECV and stacking folds, all learners and the RFECV ranker, the iterative imputer, the bootstrap resampling and the SHAP background sample (Supplementary Section~S8). Code, configuration, the fold-level results and the scripts that generate every table and figure are publicly available (see Data availability). Every number reported here was produced from repository commit \texttt{e23102d}.
